\begin{table}[H]
\caption{Síntese dos conceitos de Engenharia de Software aplicados ao projeto}
\label{tab:conceitos-engenharia-software}
\centering
\begin{tabularx}{\textwidth}{l>{\raggedright\arraybackslash}X>{\raggedright\arraybackslash}X}
\hline
\textbf{Subárea} & \textbf{Conceitos e práticas centrais} & \textbf{Aplicação direta no TCC} \\
\hline
Engenharia de Requisitos e \textit{Lean Inception} & Requisitos Funcionais (RF), Requisitos Não Funcionais (RNF), \textit{Privacy by Design}, \textit{Lean Inception}, MVP e prototipação de alta fidelidade \cite{sommerville_engenharia_2019, cavoukian_privacy_2009, ries_lean_2011, caroli_lean_2018} & Alinhamento de visão, personas e jornadas; priorização das funcionalidades do MVP de diagnóstico conceitual; conformidade \textit{LGPD-first} e validação visual das interfaces \\
\hline
Metodologias de Desenvolvimento & Desenvolvimento iterativo e incremental, Manifesto Ágil, Scrum (\textit{Product Backlog} e \textit{Sprints}) \cite{beck_manifesto_2001, schwaber_scrum_2020, cohn_user_2004} & Desdobramento do Sequenciador da \textit{Lean Inception} em \textit{Backlog} e estruturação do cronograma incremental de desenvolvimento em \textit{Sprints} \\
\hline
Design e Arquitetura de Software & Sistemas distribuídos, integração via contratos de API, referências lógicas e \textit{pipelines} intensivos em dados \cite{bass_software_2021, richards_fundamentals_2020, kleppmann_designing_2017} & Reconciliação de múltiplas fontes do Ecossistema CEDIS (Slides, Deck Manager, Atlas e Learning Passport) em um \textit{pipeline} de atributos para \textit{Knowledge Tracing} \\
\hline
\end{tabularx}
\legend{Fonte: Elaborado pelos autores.}
\end{table}

